\teachslide{Could two bases have different lengths?}{1}{
A basis gives unique coordinates.

But before defining dimension, we must prove that its length
does not depend on the chosen basis.
\spacer
Our key comparison:
\[
\key{\text{independent list length}\ \le\ \text{spanning list length}.}
\]
}{
Make the logical issue explicit. Existence and uniqueness of
coordinates refer to one chosen basis. They do not yet compare two
different bases. The following argument works over any field,
including infinite fields; it does not use the count $p^d$.
Only finite lists are involved in the comparison theorem.
}

\teachslide{Exchange one vector}{2}{
Suppose $w_1,\ldots,w_m$ span $V$ and $v\ne0$.
\[
v=a_1w_1+\cdots+a_mw_m.
\]
Some coefficient is nonzero; relabel so $a_1\ne0$.
\[
\key{w_1=a_1^{-1}\bigl(v-a_2w_2-\cdots-a_mw_m\bigr).}
\]
Thus $v,w_2,\ldots,w_m$ still span $V$.
}{
Explain both inclusions of spans. Each vector in the new list lies
in the old span. The solved equation puts the discarded $w_1$ in
the new span, and all other old vectors remain. Hence the spans are
equal. The spanning list need not be independent. Relabeling just
moves a nonzero coefficient to the first position. Dividing by it
is valid because the scalars form a field.
}

\teachslide{An exchange in $\F_3^2$}{2}{
Start with $e_1=(1,0)$ and $e_2=(0,1)$.

Insert $u=(1,1)=e_1+e_2$.
\[
e_1=u-e_2,\qquad \key{u,e_2\text{ still span}.}
\]
Now insert $v=(1,2)=u+e_2$.
\[
e_2=v-u,\qquad \key{u,v\text{ still span}.}
\]
}{
Ask students to supply each subtraction. Over $\F_3$ a minus sign
means adding twice that vector. Each exchange preserves the number
of vectors and the span. The general proof will also need to ensure
we can always remove an old $w$, rather than a vector already inserted
from the independent list. This example anticipates the retained
basis exercise later in the lecture.
}

\teachslide{Repeat the exchange}{3}{
Let $v_1,\ldots,v_r$ be independent and $w_1,\ldots,w_m$ span $V$.

After inserting $v_1,\ldots,v_k$, the list still spans and has
\[
v_1,\ldots,v_k,\quad m-k\text{ remaining old vectors}.
\]
Express $v_{k+1}$ in this list.

\key{Some remaining old vector has a nonzero coefficient.}

Otherwise $v_{k+1}\in\Span(v_1,\ldots,v_k)$, contradicting independence.
}{
Start with $k=0$. Since the list spans, such an expression exists.
If all coefficients of the remaining old vectors vanished, moving
$v_{k+1}$ to the other side would give a relation among the $v_i$
with its coefficient -1, contradicting independence. Therefore
solve for an old vector having nonzero coefficient and replace that
vector by $v_{k+1}$. The earlier exchange formula preserves spanning
and never removes an already inserted $v_i$.
}

\teachslide{An independent list cannot be longer}{2}{
Every insertion replaces one remaining old vector.

After $m$ insertions, the spanning list would be
\[
v_1,\ldots,v_m.
\]
If $r>m$, then $v_{m+1}$ would be in their span, contradicting independence.
\[
\key{r\le m.}
\]
This is the \key{Steinitz exchange argument}.
}{
The contradiction is the rigorous reason the process cannot outlast
the original spanning list. If fewer than $m$ independent vectors
were given, the process stops earlier and the inequality already
holds. If $m=0$, the space is $\{0\}$ and any nonempty independent
list is impossible, so the assertion remains true. The argument
also shows an independent finite list can be extended to a basis
when a finite spanning list is available: keep the inserted vectors
and remove any remaining redundancies.
}

\teachslide{Dimension is well defined}{2}{
Suppose $B$ and $C$ are bases, of lengths $r$ and $m$.
\begin{itemize}
\item $B$ independent, $C$ spanning: $r\le m$.
\item $C$ independent, $B$ spanning: $m\le r$.
\end{itemize}
Therefore \key{$r=m$}.

For a space with a finite basis, this common length is its dimension:
\[
\dim_k V=r.
\]
}{
Only now define dimension. Such a space is called finite-dimensional.
The result proves independence from the choice of basis, not
independence from the choice of scalar field. The theorem compares
finite bases. It also rules out an infinite independent basis when
a finite spanning list exists, since an arbitrarily long finite
sublist would violate the inequality. No Gaussian elimination was
needed, though elimination supplies another proof of the same
comparison theorem. The zero space has the empty basis and dimension
zero. The next slide specializes to finite-field coordinates.
}
